\section{Notations}
Let $\Gamma$ be a finite graph with adjacency matrix $A_\Gamma$. 
For $u, v$ lying in the vertex set of~$\Gamma$, the $(u,v)$-th entry of $A_\Gamma$ is equal to the number of edges from $u$ to $v$ if $u\neq v$, and  for $u = v$, the $(u,v)$-th entry of $A_\Gamma$ is equal to the number of loops at $u$. The matrices $T_\Gamma, T_\Gamma^\mo, L_\Gamma, \calL_\Gamma$ are defined as follows. The matrix $T_\Gamma$ is diagonal, and its $(v, v)$-th entry is equal to the degree of $v$. The matrix obtained by inverting the nonzero diagonal entries of $T_\Gamma$ is denoted by $T_\Gamma^\mo$. The matrix $T_\Gamma - A_\Gamma$ is denoted by $L_\Gamma$. The \textit{normalised Laplacian} of $\Gamma$ is denoted by $\calL_\Gamma$ and it is defined to be $T_\Gamma^{-\frac 12} L_\Gamma T_\Gamma^{- \frac 12}$. The \textit{normalised adjacency matrix} of $\Gamma$ is $T_\Gamma^{-\frac 12} A_\Gamma T_\Gamma^{- \frac 12}$. The \textit{square-graph} of $\Gamma$ is denoted by $\Gamma^2$, and it has the same set of vertices as that of $\Gamma$ and has adjacency matrix equal to $A_\Gamma^2$. The \textit{vertex-Cheeger constant} of $\Gamma$ is defined as 
$$\min_{X\subseteq V(\Gamma), 0 < |X| \leq \frac{|V(\Gamma)|}2}
\frac{|\{v\in V(\Gamma) \setminus X\,|\, v \text{ is adjacent to some element of } X\}|}{|X|}
,$$
where $V(\Gamma)$ denotes the set of vertices of $\Gamma$. 

In the following, $(V, E)$ denotes a finite graph (which may contain multiple edges, and even multiple loops at certain vertices) with $|V| \geq 4$. The maximum (resp. minimum) degree of $(V, E)$ is denoted by $d$ (resp. $d'$). The adjacency operator of $(V, E)$ is denoted by $A$. The vertex-Cheeger constant of $(V, E)$ (resp. $(V, E)^2$) is denoted by~$h$ (resp. $\tth$). Let $\scrA$ be a subset of $V$ where the vertex-Cheeger constant of $(V, E)^2$ is attained. The neighbourhood of a subset $V'$ of $V$ in $(V, E)$ is denoted by $\calN(V')$. Henceforth, we assume that Condition~\ref{Cond:MaxDegDNonBipar} holds. 
\begin{condition}
\label{Cond:MaxDegDNonBipar}
\quad 
\begin{enumerate}
\item 
The graph $(V, E)$ is undirected, that is, its adjacency matrix is symmetric. 
\item The graph $(V, E)$ is non-bipartite. 
\item The vertex-Cheeger constant $h$ of $(V, E)$ satisfies the inequality $h\geq  \varepsilon$ where $\varepsilon$ is a positive real  number. 
\end{enumerate}
\end{condition}
Set
$$
\nu = \nu_d
= 
\begin{cases}
2 & \text{ if } d = 2, \\
\frac 52 & \text{ if } d\geq 3.
\end{cases}
$$
Since $|V|\geq 4$ holds and the graph $(V, E)$ is an $\varepsilon$-vertex expander with $\varepsilon>0$, it follows that there are two distinct vertices $u, v$ of $V$ which are adjacent, and 
$$
\varepsilon |\{u, v\}|
 \leq |(\calN(u) \cup \calN(v) ) \setminus \{u, v\}| 
 \leq |\calN(u) \setminus \{v\}| + |\calN(v)\setminus \{u\}|
 \leq 2(d-1),
$$
which implies $\varepsilon \leq d-1$, and hence $\varepsilon \leq \nu -1$ if $d = 2$. If $|V|$ is odd (resp. even), then the vertex-Cheeger constant $h$ of $(V, E)$ satisfies the inequality $h \leq \frac 32$ (resp. $h \leq 1$) and hence $\varepsilon \leq \nu -1$ if $d\geq 3$. 

In the results below, we will work under Condition~\ref{Cond:SuubgraphOfRegularGraph}. 

\begin{condition}
\label{Cond:SuubgraphOfRegularGraph}
There exists an undirected graph $(V, \widetilde E)$ of degree $d$ (that is, its adjacency matrix is symmetric and has the constant vector $[1, \ldots, 1]$ as an eigenvector with eigenvalue $d$) such that $(V, E)$ is a subgraph of $(V, \widetilde E)$. 
\end{condition}

Under the above condition, the size of the difference of the edge sets of $(V, E)$ and $(V, \widetilde E)$ is denoted by $r$, the size of the difference of the edge sets of their square graphs is denoted by $s$, the vertex-Cheeger constant of $(V, \widetilde E)$ is denoted by $\widetilde h$. Note that $2r$ (resp. $2s$) is equal to the difference of the sum of the entries of the adjacency operator of $(V, \widetilde E)$ (resp. $(V, \widetilde E)^2$) and the sum of the entries of the adjacency operator of $(V, E)$ (resp. $(V, E)^2$). 

\begin{remark}
\label{Rk:NotationsCalN}
If Condition~\ref{Cond:SuubgraphOfRegularGraph} holds, then the neighbourhood of a subset $V'$ of $V$ in $(V, \widetilde E)$ is denoted by $\widetilde \calN(V')$. By the Birkhoff-von Neumann theorem~\cite[Theorem~5.5]{vanLintWilsonCourseCombi}, there exist permutations $\theta_1, \ldots, \theta_d: V\to V$ such that the vertices $v, \theta_i(v)$ are adjacent in $(V, \widetilde E)$ for any $v\in V$ and $1\leq i \leq d$, and that $\widetilde \calN(v)$ is equal to $\cup_{i=1}^d \{\theta_i(v)\}$ for any $v\in V$. For a subset $V'$ of $V$, denote the subset $\theta_i(V')$ of $\widetilde \calN(V')$ by $\widetilde \calN^i(V')$, and the subset $\theta_i(V') \cap \calN (V')$ of $\calN(V')$ by $\calN^i(V')$.
\end{remark}
